% Generated draft: requires mathematical and source review.
\documentclass[11pt]{article}
\usepackage[margin=25mm]{geometry}
\usepackage{fontspec}
\setmainfont{Latin Modern Roman}
\usepackage{amsmath,amssymb,mathtools,unicode-math}
\setmathfont{Latin Modern Math}
\usepackage{enumitem,needspace,xurl,hyperref}
\setlength{\parindent}{0pt}
\setlength{\parskip}{5pt}
\setlength{\emergencystretch}{3em}
\newcommand{\N}{\mathbb{N}}
\newcommand{\Z}{\mathbb{Z}}
\newcommand{\Q}{\mathbb{Q}}
\newcommand{\R}{\mathbb{R}}
\newcommand{\C}{\mathbb{C}}
\newcommand{\F}{\mathbb{F}}
\newcommand{\A}{\mathbb{A}}
\newcommand{\PP}{\mathbb{P}}
\newcommand{\E}{\mathbb{E}}
\newcommand{\eps}{\varepsilon}
\newcommand{\dd}{\,\mathrm{d}}
\begin{document}
\section*{UnsolvedMath ID: 30003561}
\textbf{Source problem code:} catalogue entry
\subsection{Definitions and mathematical statement}
For \(\ell\in\{2,3\}\), the supplied record distinguishes two families of arithmetic invariants: the mod-\(\ell\) Milnor invariants \(\mu_\ell(I)\) of primes, and the \(\ell\)-adic Milnor invariants \(\mu^{(\ell)}(g;I)\) of Galois elements \(g\). Here \(I\) is the index appearing in the source notation. The record does not specify its admissible entries or lengths, the underlying primes and number field, the Galois group containing \(g\), or the constructions, coefficient objects, choices, and possible indeterminacies of these invariants.

\textbf{Source question.} For each \(\ell\in\{2,3\}\), determine whether there is a relation between
\[
\mu_\ell(I)\qquad\text{and}\qquad\mu^{(\ell)}(g;I).
\]
The question imposes no explicit restriction on the length of \(I\). It does not specify which Galois elements should correspond to which prime data, what form a relation should take, or the hypotheses under which it should hold. Consequently, this is a statement of the supplied research question, not a fully specified quantified conjecture. In particular, a reduction formula or equality cannot be asserted from the record alone.

\textbf{Supplied qualification.} The dated curation labels the problem partially solved and reports that a relation is proved in quadratic and cubic cases with index length three, while the general relation remains unresolved. The record does not state those results or their hypotheses. This assessment is preserved as supplied curation, not as an independently verified mathematical fact.
\subsection{Short English statement}
For \(\ell=2\) and \(\ell=3\), ask how mod-\(\ell\) Milnor invariants attached to primes relate to \(\ell\)-adic Milnor invariants attached to Galois elements. The source leaves the underlying arithmetic data, admissible indices, correspondence, and intended relation unspecified. Its dated assessment reports partial results for indices of length three; a precise general conjecture requires further source clarification.
\subsection{Sources}
\begin{itemize}
\item Local UnsolvedMath registry, supplied complete record: permanent ID \(30003561\), title \emph{Relations Between Arithmetic Milnor Invariants}, exact source problem code \texttt{catalogue entry}; collection: Oberwolfach Reports Open Problems. No individual proposer is supplied. The proposed year is recorded as 2017.
\item Supplied report attribution: \emph{Low-dimensional Topology and Number Theory} (2018), Oberwolfach Reports, pp.~2363--2426, DOI \texttt{10.4171/owr/2017/38}. \url{https://doi.org/10.4171/owr/2017/38}. The background describes this collection as drawn from mathematical workshop reports and their open-problem discussions.
\item Supplied literature reference: \emph{Arithmetic topology in Ihara theory II}. \url{https://doi.org/10.1016/j.jnt.2018.10.010}. No theorem number, authors, or precise comparison statement is supplied.
\item Supplied literature reference: \emph{Low dimensional Topology and Number Theory}. \url{https://doi.org/10.4171/owr/2017/38}.
\item Literature assessment date: \texttt{2026-08-21}. This is supplied dated curation, not independently verified. It labels the problem partially solved, reports a proved relation in quadratic and cubic cases with index length three, and describes the unrestricted general relation as unresolved. The supplied background explicitly requests independent verification of time-sensitive status and cited claims.
\item The registry labels the statement \texttt{corrected\_verified} and says that mathematical content was preserved while notation, typography, or omitted definitions were repaired. This is a supplied registry qualification, not independent verification of the formulation or its mathematical correctness.
\end{itemize}
\end{document}
